Time series imputation in real applications is often performed under mixed missingness, where MAR-like and MNAR-like patterns co-exist. Prior methods usually rely on a shared backbone that mainly emphasizes observed-to-missing recovery, while missing-to-missing dependency is only modeled implicitly. This mismatch becomes critical when different missing positions require different recovery evidence.

We present UniMiss, a mechanism-aware imputation framework that explicitly decomposes recovery into four components: an \textbf{OO foundational context path} for stable observed-side representation, an \textbf{OM interaction-oriented expert branch} for observed evidence routing, an \textbf{MM missing-structure expert branch} for co-missing topology and extreme/periodic clues, and a \textbf{Stage-II MAR-vs-MNAR soft gate} that adaptively balances OM and MM at each missing position. The design follows the local CIKM draft objective of separating ``who provides stable context,'' ``who models missing dependencies,'' and ``how mechanism preference is selected under mixed settings.''

Our experimental protocol follows the unified benchmark setting in this project: ETT and IAQ datasets, MAR/MNAR/Mix mechanisms, missing rates of 20\%, 30\%, and 40\%, and multi-seed evaluation with MAE/RMSE/MRE/NRMSE. We provide switchable module-level ablations, hyperparameter studies, visualization diagnostics, and lightweight/scalable variants within one reproducible code framework. The resulting manuscript and repository are organized to support full main-table, ablation, and analysis reporting under a single mechanism-aware protocol.
